\chapter{Bias and fairness check BIAS-1}
\label{ch:bias}

The governance agent runs a bias and fairness check, BIAS-1, on the resident consultation responses. It looks for skew by subgroup and by data-collection channel, the kind of bias assessment that AI risk management guidance asks for \parencite{nist2023airmf}.

\section{Method}

The method is implemented in \code{app/governance.py:43-107}:
\begin{itemize}
\item Tag each of the 214 consultation responses in \code{05_Resident_consultation_responses.csv} with themes, using a lexicon in English, Chinese, Malay and Tamil.
\item Rank the themes three ways: by raw count, with each age band weighted equally, and with English-only responses.
\item Report skew by channel, age and language.
\end{itemize}

\section{Findings}

Table~\ref{tab:bias} lists the six findings recorded in \code{docs/governance/evidence/bias_check.json}. Figure~\ref{fig:bias} draws the response counts by age band and channel, and the raw and age-re-weighted ranks of the top four themes, from the same file.

\begin{table}[htbp]
\centering
\small
\caption{Bias and fairness findings of BIAS-1}
\label{tab:bias}
\begin{tabularx}{\textwidth}{>{\raggedright\arraybackslash}p{0.7cm}>{\raggedright\arraybackslash}p{1.5cm}>{\raggedright\arraybackslash}X}
\toprule
\textbf{ID} & \textbf{Severity} & \textbf{Finding} \\
\midrule
B1 & \textcolor{pmOrange}{\textbf{high}} & Online channel gave 160 of 214 responses, 74.8\%, and 115 of them came from under-45s. Seabreeze saw the same skew, LL-L06. \\
B2 & \textcolor{pmOrange}{\textbf{high}} & Residents 65+ gave 25 responses, 11.7\%, although they are the main users of the active ageing centre; under-45s gave 121, 56.5\%. The estate's age profile is NOT FOUND in the pack, so true representativeness cannot be computed. \\
B3 & \textcolor{pmOrange}{\textbf{high}} & All 19 non-English responses, zh 10, ms 6 and ta 3, came on paper or at the town hall, 18 of them from residents 45+. An English-only summary would drop 11 of 36 active-ageing comments and 6 of 37 lift/access comments. \\
B4 & medium & Raw top four: Childcare, Communication \& trust, Lifts \& access, Active ageing. Age-re-weighted top four: Childcare, \textbf{Active ageing}, Lifts \& access, Communication \& trust. The order changes, so both must be reported. \\
B5 & medium & Only 25 distinct comment texts across 214 responses; check for template or campaign submissions before weighting. \\
B6 & low & Mean rating by age band ranges from 2.65 for 35--44 to 3.48 for 18--24. The town hall drew about 40 residents, EDC-4.5, and gave 31 responses. \\
\bottomrule
\end{tabularx}
\end{table}

% Figure: bias check charts (source: bias_check.json)
\begin{figure}[htbp]
\centering
\begin{minipage}[t]{0.56\textwidth}
\centering
\begin{tikzpicture}
\begin{axis}[
  ybar, width=\linewidth, height=5.0cm, bar width=13pt,
  ymin=0, ymax=68,
  symbolic x coords={18--24,25--34,35--44,45--54,55--64,65+},
  xtick=data,
  xticklabel style={font=\footnotesize},
  yticklabel style={font=\footnotesize},
  ylabel={responses}, ylabel style={font=\footnotesize},
  nodes near coords, nodes near coords style={font=\footnotesize},
  axis x line*=bottom, axis y line*=left,
  enlarge x limits=0.1,
  title={a. Responses by age band, $n=214$}, title style={font=\small},
]
\addplot[bar shift=0pt, fill=pmBlue, draw=pmBlue] coordinates {(18--24,21) (25--34,57) (35--44,43)};
\addplot[bar shift=0pt, fill=pmNavy, draw=pmNavy] coordinates {(45--54,31) (55--64,37)};
\addplot[bar shift=0pt, fill=pmOrange, draw=pmOrange] coordinates {(65+,25)};
\end{axis}
\end{tikzpicture}
\end{minipage}\hfill
\begin{minipage}[t]{0.40\textwidth}
\centering
\begin{tikzpicture}
\begin{axis}[
  ybar, width=\linewidth, height=5.0cm, bar width=16pt,
  ymin=0, ymax=185,
  symbolic x coords={online,paper,town hall},
  xtick=data,
  xticklabel style={font=\footnotesize},
  yticklabel style={font=\footnotesize},
  ylabel={responses}, ylabel style={font=\footnotesize},
  nodes near coords, nodes near coords style={font=\footnotesize},
  axis x line*=bottom, axis y line*=left,
  enlarge x limits=0.25,
  title={b. Responses by channel}, title style={font=\small},
]
\addplot[bar shift=0pt, fill=pmTeal, draw=pmTeal] coordinates {(online,160) (paper,23) (town hall,31)};
\end{axis}
\end{tikzpicture}
\end{minipage}

\vspace{0.6cm}
\begin{tikzpicture}[font=\footnotesize]
  \def\dx{4.6}
  \node[font=\small] at (\dx/2, 0.95) {c. Top four themes: raw rank vs age-re-weighted rank};
  \node[font=\footnotesize\bfseries, anchor=east] at (0, 0.35) {Raw: count};
  \node[font=\footnotesize\bfseries, anchor=west] at (\dx, 0.35) {Age-re-weighted: weight};
  % raw ranks (y = -rank * 0.65)
  \foreach \r/\t/\v in {1/Childcare/45, 2/Communication \& trust/38, 3/Lifts \& access/37, 4/Active ageing/36}{
    \node[anchor=east] (L\r) at (-0.15, -0.65*\r+0.35) {\r.\ \t: \v};
    \fill[pmNavy] (0, -0.65*\r+0.35) circle (2pt);
  }
  \foreach \r/\t/\v in {1/Childcare/0.2061, 2/Active ageing/0.1800, 3/Lifts \& access/0.1774, 4/Communication \& trust/0.1772}{
    \node[anchor=west] (R\r) at (\dx+0.15, -0.65*\r+0.35) {\r.\ \t: \v};
    \fill[pmNavy] (\dx, -0.65*\r+0.35) circle (2pt);
  }
  % connectors: raw rank -> reweighted rank
  \draw[pmGrey, line width=1pt] (0,-0.30) -- (\dx,-0.30);           % Childcare 1 -> 1
  \draw[pmGrey, line width=1pt] (0,-1.60) -- (\dx,-1.60);           % Lifts 3 -> 3
  \draw[pmOrange, line width=1.6pt] (0,-2.25) -- (\dx,-0.95);       % Active ageing 4 -> 2
  \draw[pmBlue, line width=1.6pt] (0,-0.95) -- (\dx,-2.25);         % Communication 2 -> 4
\end{tikzpicture}
\caption[Bias check: response skew and theme ranks]{Bias check BIAS-1. Panel a: responses by age band, light blue under 45, navy 45--64, orange 65+; residents 65+ gave 25 or 11.7\%, under-45s 121 or 56.5\%. Panel b: responses by channel; online gave 160 of 214, 74.8\%. Panel c: rank of the top four themes by raw count and with each age band weighted equally. Active ageing rises from fourth to second and Communication \& trust falls from second to fourth, finding B4. Re-weighted values are the theme weights in the file; theme names are shortened as in finding B4. Source: \texttt{bias\_check.json}, fields \texttt{age}, \texttt{channel}, \texttt{themes\_raw}, \texttt{themes\_reweighted}, \texttt{rank\_raw} and \texttt{rank\_reweighted}.}
\label{fig:bias}
\end{figure}


\section{Actions proposed by the check}

\begin{enumerate}
\item Translate and include every non-English comment.
\item Show raw and age-re-weighted priorities side by side.
\item Run a second multilingual paper or door-to-door round for residents 55+; see LL-L06.
\item Have Community Engagement confirm whether the repeated submissions are independent.
\item Hold the consultation summary D15 until HC7.
\end{enumerate}

\section{Status and team decision}

\paragraph*{Status.} The Supervisor account signed BIAS-1 off in the demonstration run; this is approval \#10 in Table~\ref{tab:approvals}. As noted in Section~\ref{sec:demoapprovals}, it is a demonstration decision, not the team's.

\paragraph*{What the team changed as a result: the team's HC7 decision.} The team approved and carried out the proposed actions except the most critical one, action 3, the second multilingual paper / door-to-door round for residents 55+, which was not approved. So non-English comments are translated and included, action 1; raw and age-re-weighted priorities are shown side by side, action 2; repeated submissions are checked, action 4; and the summary D15 was held until HC7, action 5. \textbf{The under-representation of residents 65+, finding B2, therefore remains} and must be stated in the consultation summary.
